% ============================================================
% SynthDetect: AI-Generated Text Detection via
% Disrupt-and-Recover Analysis and Fine-Grained Attribution
% ============================================================
\documentclass[conference]{IEEEtran}

\usepackage{amsmath,amssymb,amsfonts}
\usepackage{algorithmic}
\usepackage{graphicx}
\usepackage{textcomp}
\usepackage{xcolor}
\usepackage{hyperref}
\usepackage{booktabs}
\usepackage{multirow}
\usepackage{subcaption}

\hypersetup{
    colorlinks=true,
    linkcolor=blue,
    citecolor=blue,
    urlcolor=blue,
}

\begin{document}

\title{SynthDetect: AI-Generated Text Detection via\\
Disrupt-and-Recover Analysis and Fine-Grained Attribution}

\author{
\IEEEauthorblockN{SynthDetect Team}
\IEEEauthorblockA{%
\textit{Research Paper — Work in Progress}\\
2026}
}

\maketitle

% ============================================================
\begin{abstract}
We present SynthDetect, a novel framework for detecting AI-generated text
that combines two complementary approaches: (1) \textbf{Disrupt-and-Recover (D\&R)},
which exploits posterior concentration in language models by measuring how
faithfully an LLM reconstructs deliberately shuffled text, and
(2) \textbf{Fine-Grained AI Detection (FAID)}, which uses contrastive learning
on multi-level linguistic features for model-family attribution.
Our key insight is that AI-generated text lies in high-probability regions
of a model's output distribution, causing disrupted AI text to be recovered
with significantly higher fidelity than human-written text.
The D\&R pipeline requires only black-box API access (no log-probabilities),
making it robust against model updates.
Preliminary experiments using Google Gemini 2.5 Flash show a recovery
similarity gap of $\Delta = +0.0507$ between AI and human text,
with structural similarity providing the strongest discriminative signal
($+0.1243$). SynthDetect achieves this with a single LLM call per text
chunk at a cost of less than \$0.001 per document.
\end{abstract}

\begin{IEEEkeywords}
AI-generated text detection, language model attribution, contrastive learning,
posterior concentration, natural language processing
\end{IEEEkeywords}

% ============================================================
\section{Introduction}

The rapid proliferation of large language models (LLMs) such as GPT-4,
Claude, and Gemini has created an urgent need for reliable methods to
distinguish AI-generated text from human writing. Existing detection
approaches fall into several categories: watermarking-based methods
\cite{kirchenbauer2023}, which require model cooperation; statistical
methods based on log-probability analysis \cite{mitchell2023}, which
require white-box access; and classifier-based methods that may not
generalize across model families \cite{guo2023}.

We propose SynthDetect, a framework that addresses these limitations
through two novel contributions:

\begin{enumerate}
    \item \textbf{Disrupt-and-Recover (D\&R)}: A black-box detection method
    that exploits the \textit{posterior concentration} phenomenon in LLMs.
    We deliberately disrupt text structure through controlled sentence
    shuffling, then ask an LLM to recover the original order. AI-generated
    text, lying in high-probability regions of the model's output
    distribution, is recovered with significantly higher fidelity.

    \item \textbf{Fine-Grained AI Detection (FAID)}: A contrastive learning
    approach that maps multi-level linguistic features (lexical, syntactic,
    semantic, stylometric) into an embedding space where texts from the same
    source cluster together, enabling model-family attribution.
\end{enumerate}

% ============================================================
\section{Related Work}

\subsection{Statistical Detection Methods}
Mitchell et al.\ \cite{mitchell2023} proposed DetectGPT, which uses
log-probability curvature to identify machine-generated text. While
effective, this approach requires white-box access to model probabilities.

\subsection{Classifier-Based Approaches}
OpenAI's text classifier \cite{solaiman2019} and similar supervised
approaches train binary classifiers on human vs.\ AI text. These methods
often suffer from distribution shift when new models emerge.

\subsection{Watermarking}
Kirchenbauer et al.\ \cite{kirchenbauer2023} introduced soft watermarking
for LLM outputs. While robust when applicable, watermarking requires
cooperation from model providers and does not address existing unwatermarked text.

% ============================================================
\section{Methodology}

\subsection{Disrupt-and-Recover (D\&R) Pipeline}

The D\&R pipeline consists of four stages:

\subsubsection{Text Chunking}
Input text $T$ is segmented into chunks $\{C_1, C_2, \ldots, C_n\}$
using semantic chunking that respects sentence boundaries, with a
target size of $\sim$200 words per chunk.

\subsubsection{Controlled Disruption}
Each chunk $C_i$ undergoes sentence-level shuffling $\sigma(C_i)$,
producing a disrupted version $\tilde{C}_i$. We preserve the first and
last sentences as context anchors and maintain 20\% of middle sentences
in their original positions.

\subsubsection{LLM Recovery}
The disrupted chunk $\tilde{C}_i$ is submitted to an LLM (Google Gemini
2.5 Flash) with a deterministic recovery prompt ($T=0.0$), producing
a recovered version $\hat{C}_i$.

\subsubsection{Similarity Scoring}
We compute a hybrid similarity score:
\begin{equation}
    S(C_i, \hat{C}_i) = \alpha \cdot S_{\text{sem}}(C_i, \hat{C}_i)
    + (1 - \alpha) \cdot S_{\text{str}}(C_i, \hat{C}_i)
\end{equation}
where $S_{\text{sem}}$ is the cosine similarity of sentence-transformer
embeddings (all-MiniLM-L6-v2), $S_{\text{str}}$ is the normalised
Kendall-$\tau$ distance of sentence orderings, and $\alpha = 0.6$.

The final D\&R score aggregates across chunks:
\begin{equation}
    S_{DR}(T) = \frac{1}{n} \sum_{i=1}^{n} S(C_i, \hat{C}_i)
\end{equation}

\subsection{Fine-Grained AI Detection (FAID)}

FAID extracts multi-level features (lexical: 10 dims, syntactic:
$\sim$47 dims, semantic: 4 dims, stylometric: 15 dims) and maps them
through a contrastive encoder trained with SupConLoss
\cite{khosla2020}. k-NN search in the resulting embedding space
provides model-family attribution.

\subsection{Fusion Layer}

The fusion layer combines both signals:
\begin{equation}
    S_{\text{final}} = w_{DR} \cdot S_{DR} + w_{FAID} \cdot S_{FAID}
\end{equation}
with $w_{DR} = 0.6$ and $w_{FAID} = 0.4$, along with variance-based
collaborative text detection.

% ============================================================
\section{Preliminary Results}

\subsection{Prototype Validation}

Initial experiments on sample texts using Gemini 2.5 Flash demonstrate
the core hypothesis (Table~\ref{tab:results}).

\begin{table}[h]
    \centering
    \caption{D\&R Pipeline — Preliminary Results}
    \label{tab:results}
    \begin{tabular}{lcc}
        \toprule
        \textbf{Metric} & \textbf{Human} & \textbf{AI} \\
        \midrule
        Mean Combined Sim & 0.9199 & 0.9706 \\
        Std Combined Sim & 0.0801 & 0.0294 \\
        Mean Semantic Sim & 0.9937 & 0.9953 \\
        Mean Structural Sim & 0.8091 & 0.9334 \\
        \midrule
        $\Delta$ (AI $-$ Human) & \multicolumn{2}{c}{\textbf{+0.0507}} \\
        \bottomrule
    \end{tabular}
\end{table}

Key observations:
\begin{itemize}
    \item Structural similarity shows the strongest discriminative signal
    ($\Delta = +0.1243$), indicating that sentence order recovery is the
    primary differentiator.
    \item AI text exhibits lower variance across chunks ($\sigma^2 = 0.0294$
    vs.\ $0.0801$), consistent with posterior concentration.
    \item Total cost per document: $<$\$0.001 (4 API calls).
\end{itemize}

% ============================================================
\section{Conclusion and Future Work}

We have demonstrated that the Disrupt-and-Recover approach provides a
viable black-box method for AI text detection. Future work includes:
(1) large-scale evaluation on HC3 and M4GT-Bench datasets,
(2) training the FAID contrastive encoder on real feature data,
(3) multi-model cross-validation (GPT-4, Claude, Gemini),
(4) adversarial robustness evaluation, and
(5) threshold optimisation via ROC analysis.

% ============================================================
\begin{thebibliography}{00}

\bibitem{mitchell2023}
E. Mitchell, Y. Lee, A. Khazatsky, C. D. Manning, and C. Finn,
``DetectGPT: Zero-shot machine-generated text detection using probability
curvature,'' in \textit{Proc. ICML}, 2023.

\bibitem{kirchenbauer2023}
J. Kirchenbauer, J. Geiping, Y. Wen, J. Katz, I. Miers, and T. Goldstein,
``A watermark for large language models,'' in \textit{Proc. ICML}, 2023.

\bibitem{guo2023}
B. Guo, X. Zhang, Z. Wang, M. Jiang, et al.,
``How close is ChatGPT to human experts? Comparison corpus, evaluation,
and detection,'' 2023, arXiv:2301.07597.

\bibitem{solaiman2019}
I. Solaiman, M. Brundage, J. Clark, et al.,
``Release strategies and the social impacts of language models,'' 2019,
arXiv:1908.09203.

\bibitem{khosla2020}
P. Khosla, P. Teterwak, C. Wang, A. Sarna, Y. Tian, P. Isola, A. Maschinot,
C. Liu, and D. Krishnan, ``Supervised contrastive learning,'' in
\textit{Proc. NeurIPS}, 2020.

\end{thebibliography}

\end{document}
